\section{Limitations}\label{sec:limitations}

This report measures one server pair, one database, one schema, and one scenario set. Server versions matter (SQLcl 26.1, OraViz 0.1.0, Oracle AI Database 26ai Free 23.26), formatting differs across SQLcl releases, and the demo data is small by design so that the framing overhead of small results remains visible. The per-stage ratios should therefore be read as configuration-specific measurements; the structural properties (boundedness of the contract, linear growth of unbounded results, fixed schema cost) are what generalize.

Token counts use \texttt{tiktoken} as a proxy for the tokenizer of any particular agent model, and images are excluded from the accounting entirely; a model that bills images as tokens would see the chart path differently, and a model with a very different vocabulary might shift the small-result comparison. The two encodings we report agree within one point on the aggregate, which bounds but does not eliminate this concern. The measurements are a single deterministic run: the database is static and the outputs are text, so repeated runs reproduce the same numbers on the same versions, but they do not capture network latency, retries, or agent-level variance.

The evaluation measures context cost, not task success. A smaller context is only better if the agent can still answer correctly, and the contract relies on explicit truncation markers for that. We did not run a controlled accuracy study comparing agents that receive bounded results against agents that receive full ones; for questions whose evidence lies beyond the preview, the bounded server requires a follow-up call, and whether that exchange is net-positive is an empirical question this report does not answer. Relatedly, the chart path is validated end to end through an agent loop, but the quality of the resulting chart for a given question is the agent's responsibility.

Finally, the safety posture is deliberately modest. The read-only guarantee is lexical validation, not a database-enforced boundary, and the network transports have no built-in authentication, so deployments that expose HTTP beyond localhost should add a proxy or a read-only database account. These choices are documented in the artifact rather than presented as solved, because the report's subject is context efficiency, and overstating either property would undercut the measurements it presents.
